% Une période du protocole SWIM : ping, ping-req par des intermédiaires, puis suspicion.
\documentclass[border=6pt]{standalone}
\usepackage{sd}
\begin{document}
\begin{tikzpicture}[x=1mm, y=1mm, note/.append style={font=\sffamily\small}]
  \foreach \x/\t in {0/{membre i}, 38/{intermédiaires m1..m3}, 76/{membre j}} {
    \node[pillo=Dslate] at (\x,70) {\t};
    \draw[seqlife=Dslate] (\x,66) -- (\x,0);
  }
  \draw[arr=Dblue] (0,62) -- node[lblc=Dblue, above, pos=0.3] {1. ping} (60,57);
  \node[note, text=Dred, anchor=west] at (61,56) {perdu};
  \node[note, anchor=east, align=right] at (-2,52) {pas d'ack\\avant le délai};
  \draw[arr=Dblue] (0,46) -- node[lblc=Dblue, above] {2. ping-req(j), à k = 3} (38,42);
  \draw[arr=Dblue] (38,40) -- node[lblc=Dblue, above] {3. ping} (76,36);
  \draw[arr=Dgreen] (76,32) -- node[lblc=Dgreen, above] {ack} (38,28);
  \draw[arr=Dgreen] (38,26) -- node[lblc=Dgreen, above] {4. ack relayé} (0,22);
  \node[note, anchor=west, align=left] at (78,34) {j répond : vivant,\\le ping direct était\\seulement perdu};
  \node[note, anchor=north, align=center, fill=white] at (38,16) {sans aucune réponse : i déclare j suspect, la nouvelle voyage sur les messages ;\\j réfute s'il est vivant (incarnation + 1), sinon il est déclaré mort\\après 4 × log₁₀ n périodes};
\end{tikzpicture}
\end{document}
